\section{CLASP: The Paired Cross-Layer Audit}
\label{sec:framework}

\subsection{Overview}
\label{sec:framework_overview}

CLASP compares two representations of the same observations. Within each
country--measurement unit, every candidate-bearing segment contributes one
unit of mass to a network-transition distribution and one unit, allocated
among candidates, to a feasible-corridor distribution. Figure
~\ref{fig:framework_pipeline} places this paired construction at the center
of the workflow: path extraction defines the observations, corridor
projection supplies feasible support, and allocation and aggregation produce
the distributions used for within-unit and cross-unit comparisons.

The inputs are traceroutes, probe metadata, IP geolocation and ASN
mappings, AS context, and cable metadata. The main branch conditions both
distributions on candidate-bearing segments. A separate trace-level branch
reports exposure using all valid traceroutes and the presence of feasible
inter-region candidates. Keeping these branches separate makes their
different populations and denominators explicit.

\begin{figure*}[t]
  \centering
  \includegraphics[width=\textwidth]{figures/framework_cross_layer_audit.pdf}
  \caption{CLASP preserves the same candidate-bearing segments and total
  observation mass in the two distributions. The main flow separates inputs,
  feasible support, paired construction, and comparison. The trace-level
  exposure branch uses all valid traces and inter-region candidate flags;
  dashed arrows denote supporting evidence or flags.}
  \label{fig:framework_pipeline}
\end{figure*}

\subsection{Path Extraction}
\label{sec:framework_path_segmentation}

Country and measurement identity define the comparison unit before
segments are aggregated. We group traceroutes by probe country and measured
service, retaining source-side connectivity and observed destination
selection. Hop countries remain segment attributes rather than grouping
keys, so intermediate locations do not redefine the source population.

Each traceroute is normalized into its complete sequence of visible hops with
locations, ASNs, and RTTs. Successive visible, geolocated hops form atomic
segments. Each segment records endpoint locations, ASNs, countries, hop
positions, RTT difference, and bridged-hop count. Repeated hops and RTT samples
retain their trace identity and ordering. When timeouts separate two visible
hops, the bridged-hop count records the hidden gap for audit accounting.

When both endpoint ASNs are available, their ordered pair forms the segment's
network-transition label. A segment missing one endpoint ASN receives an
ordered country-pair fallback label, which is strictly coarser. We record the
fallback share for each unit. Audit eligibility requires that share to be at
most 30\%, together with at least 30 candidate-bearing segments, 10 probes,
and 3 probe ASNs. These thresholds remove units whose network labels or
observation support are too sparse for the paired summaries; they are quality
floors rather than guarantees of representativeness.

\subsection{Feasible Landing-Region Corridors}
\label{sec:framework_projection}

\paragraph{Landing stations and cable lifecycle.}
For each segment endpoint, CLASP retrieves landing stations within 50~km and
pairs them through cable systems in service on July~1, 2026. The catchment
allows endpoint locations to be associated with nearby coastal facilities at
the resolution of the geolocation input. Explicit cable path or branch
information defines direct landing-point adjacency. When the inventory does
not provide that information, the pipeline retains reachability at the
metadata's topology resolution and records that provenance state rather than
presenting it as an observed cable traversal.

\paragraph{Landing regions.}
Landing stations are grouped into regions with maximum pairwise distance
30~km. The diameter merges nearby facilities serving the same coastal
direction while retaining separated coastal entry and exit patterns. A
\emph{landing-region corridor} is a direction-independent pair of regions
connected by at least one feasible cable candidate. Several cable systems and
exact landing pairs may support one corridor. Exact landing-pair, cable, and
corridor identities remain linked in the outputs so that the aggregation can
be audited at each resolution.

\paragraph{Propagation feasibility.}
A candidate must satisfy geographic, lifecycle, and propagation constraints.
Let \(D\) be the great-circle distance between its landing points,
\(v_f=200\)~km/ms the configured effective propagation speed in fiber,
\(\Delta RTT\) the segment's RTT difference, and \(\tau=5\)~ms the RTT
tolerance. The candidate is propagation-feasible when

\begin{equation}
  \Delta RTT + \tau \geq \frac{2D}{v_f}.
  \label{eq:rtt-feasibility}
\end{equation}

The inequality asks whether the observed round-trip increment can accommodate
the minimum round-trip propagation time between the landing points. When
\(\Delta RTT\) is non-positive or inconclusive, CLASP retains the geographic
and lifecycle constraints and flags the segment. The 200~km/ms speed and
5~ms tolerance are fixed configuration values in the reported analysis.

\subsection{Candidate Allocation and Observation Mass}
\label{sec:framework_allocation}

Candidate multiplicity must not increase a segment's observation mass.
Duplicate cable rows belonging to the same corridor are collapsed within
an atomic segment. For a segment \(s\) with feasible corridor set
\(\mathcal{C}_s\), projection-score allocation distributes one unit of mass
according to the candidates' relative support:

\begin{equation}
  w_{s,c}=
  \frac{\sigma_{s,c}}
       {\sum_{c'\in\mathcal{C}_s}\sigma_{s,c'}},
  \qquad
  \sum_{c\in\mathcal{C}_s}w_{s,c}=1,
  \label{eq:weighted-mass}
\end{equation}

where \(\sigma_{s,c}\) sums the evidence-support scores of the cable candidates
assigned to corridor \(c\). The implementation forms each candidate score
as the product of landing-proximity support, AS-economic support, RTT
support, and a latency penalty. Landing-proximity support decreases with
the distances from the two hop endpoints to their candidate landing points.
AS-economic support uses endpoint-AS membership in cable-owner AS sets and
relationships or owner-group reachability in the AS graph. RTT support
compares the observed increment with the propagation lower bound and applies
a penalty to highly inflated short-path increments; inconclusive increments
receive neutral RTT factors. The candidate scores are summed by corridor
and normalized as in Eq.~\eqref{eq:weighted-mass}, so every candidate-bearing
segment contributes one unit of mass.

The three allocation rules retain different amounts of candidate-score
information while preserving this unit-mass constraint. Projection-score
weighting is the primary allocation for Layer Agreement and retains
relative evidence support. Equal sharing assigns
\(w_{s,c}=1/|\mathcal{C}_s|\) without using score magnitudes, and Top-1
assigns all mass to the highest-scoring corridor. Endpoint-AS information
contributes both to network labels and to candidate ranking; equal sharing
therefore provides a comparison that does not use the score magnitudes
within the feasible sets. The rules express equal, graded, and single-
candidate support, not confirmed cable usage probabilities.

\subsection{Mapping Resolution and Pipeline Accounting}
\label{sec:framework_mapping_resolution}

Mapping states determine which segments enter the paired distributions.
Table~\ref{tab:mapping-resolution} distinguishes input coverage and candidate-
set size. Single- and bounded multi-corridor segments enter the physical
distribution and its paired network distribution. Segments with no feasible
corridor or insufficient inputs remain in pipeline accounting; every valid
traceroute remains in the exposure denominator.

\begin{table}[t]
  \centering
  \caption{Segment-level physical mapping resolution.}
  \label{tab:mapping-resolution}
  \footnotesize
  \setlength{\tabcolsep}{3.2pt}
  \begin{tabular}{p{0.30\columnwidth}p{0.60\columnwidth}}
    \toprule
    State & Operational definition \\
    \midrule
    Single corridor
      & Feasibility inputs available and \(|\mathcal{C}_s|=1\) \\
    Bounded multi-corridor
      & Feasibility inputs available and \(|\mathcal{C}_s|>1\) \\
    No feasible corridor
      & Feasibility inputs available and \(|\mathcal{C}_s|=0\) \\
    Insufficiently resolved
      & Insufficient location, path-visibility, or metadata fields for stable
        corridor construction \\
    \bottomrule
  \end{tabular}
\end{table}

Candidate-set size and aggregate concentration describe different levels. A
unit may contain many bounded multi-corridor segments yet concentrate after
their allocated mass repeatedly supports the same corridors. CLASP therefore
retains the resolution state instead of using candidate multiplicity as a
substitute for distribution breadth.

\subsection{Paired Distributions}
\label{sec:framework_distribution_analysis}

For an auditable country--measurement unit \(u\), let \(S_u\) denote its
candidate-bearing atomic segments and \(\ell(s)\) the ordered network label of
segment \(s\). The network-transition probability of label \(t\) is

\begin{equation}
  p_u^{N}(t)=
  \frac{\sum_{s\in S_u}\mathbf{1}[\ell(s)=t]}
       {|S_u|}.
  \label{eq:network-distribution}
\end{equation}

The corridor distribution sums the allocated masses in
Eq.~\eqref{eq:weighted-mass} and normalizes over the same \(S_u\).
Before normalization, each layer has total mass \(|S_u|\): candidate
ambiguity changes the allocation, not the amount contributed by a segment.
The paired distributions thus compare how a fixed observation population
is organized. Network labels preserve direction, whereas corridors are
direction-independent coastal connections; the comparison retains this
difference in category meaning.

Submarine exposure uses the complete valid-trace population. A traceroute is
exposed when at least one of its segments has a feasible inter-region corridor.
The number of exposed traces divided by all valid traces in the unit gives the
exposure rate. This denominator differs intentionally from \(|S_u|\), which
conditions the distributions on candidate-bearing segments.

\subsection{Distribution Summaries and Layer Agreement}
\label{sec:framework_metrics}

Within-unit comparisons measure concentration and breadth separately.
Top-2 share sums the two largest probabilities, recording mass in the
dominant categories. Effective category count,
\(\exp[-\sum_i p_i\log p_i]\), expresses breadth as an equally weighted
category count. Normalized entropy divides Shannon entropy by the logarithm
of observed support size and describes evenness relative to that support.
For each metric, network-to-corridor differences are computed within
the same unit.

\emph{Layer Agreement} is the Spearman rank association, across a reported
set of units, between network Top-2 share and corridor Top-2 share. A
coefficient approaching +1 indicates closely aligned concentration rankings;
a coefficient near zero indicates little monotonic rank association, and a
negative coefficient indicates opposing rankings. This comparison concerns
ordering across units. Within-unit Top-2 differences and effective category
counts separately measure how concentration and breadth change across
layers. A group can therefore exhibit aligned rankings alongside substantial
within-unit narrowing. We report Layer Agreement by measurement family,
geography, and candidate-allocation rule.

For the frozen equal-share result, we quantify geographic uncertainty at the
level where geography is assigned. A cluster-bootstrap replicate samples
countries with replacement within each geographic group and carries every
country--measurement unit from each sampled country together. We use 50,000
replicates and report percentile intervals. A separate two-sided reference
test permutes the geographic labels of whole countries 100,000 times while
retaining the observed number of countries in each group. Both procedures use
seed 20260919. This country-clustered analysis is applied to the frozen equal-
share unit table; projection-score and Top-1 results provide the allocation-
robustness comparison.

A common 80\% Top-2 threshold supplies four descriptive classes: broad at both
layers, concentrated at both, network-broad/corridor-concentrated, and
network-concentrated/corridor-broad. Continuous Top-2 shifts, effective counts,
and normalized entropy accompany the classes so that threshold crossings do
not replace the underlying distributions.

\begin{table}[t]
  \centering
  \caption{CLASP measurement configuration.}
  \label{tab:measurement-configuration}
  \footnotesize
  \setlength{\tabcolsep}{3.5pt}
  \begin{tabular}{lr}
    \toprule
    Parameter & Value \\
    \midrule
    Endpoint landing catchment & 50 km \\
    Landing-region maximum diameter & 30 km \\
    Fiber propagation speed & 200 km/ms \\
    RTT tolerance & 5 ms \\
    Same-city threshold & 25 km \\
    Cable lifecycle date & July 1, 2026 \\
    Inconclusive RTT & Retained and flagged \\
    Timeout-bridged segment & Retained and flagged \\
    Path scope & Complete visible sequence \\
    \bottomrule
  \end{tabular}
\end{table}
